%+PATCH,//WAVE/TEX
%+DECK,introduction,T=LATEX.

\renewcommand*\thesection{\Roman{section}}
\section{Introduction}

\W~has been developed at BESSY over the years to track electrons
through magnetic fields and to calculate spontanous synchrotron radiation
\cite{dok}.
This document illustrates the usage of \W~and makes you
familiar with it by following the given examples.

\W is mainly controlled by an ASCII input file named \wie.
For each example, there is a script \gtis{wave\_example\_n.sh} to
create \wi from \wee, where \gtis{n}~is the number of the example. Once
\wi is created by the corresponding script, you may compare
\we and \wi to see, which variables have changed.

There is a graphical user interface (GUI), a Python script called
\WSH to setup \wie, to run \W, and to visualize the
results. This is explained in the first example. However, not all
options for \W are available in the GUI.

Once \W is installed, you should have a \W directory and shell
variable \W pointing to it. It is assumed, that the user knows to edit
a file and how to open a terminal, e.g. \gtis{bash} under Linux  or powershell and under
Windows. As an alternative one may run the Python shell \gtise{idle}.

WAVE and the related software is published under the GNU General Public
License. To get WAVE, please contact the author or get it from github:

~~~~~\gtis{git clone https://github.com/MicScheer/wave.git wave}

